\documentclass{article}

\begin{document}

\section*{Qu'est-ce qu'une dérivée ?}
\subsubsection*{Les dérivées expliquées à ma petite nièce...}

\wcSpan*[class=tp-red-style]{— C’est quoi une dérivée ?}

Eh bien, c’est simple : c’est la limite du taux d’accroissement.

\wcSpan*[class=tp-red-style]{— Euh…}

Oui, le nombre dérivé d’une fonction $f$ en un point $a$, c’est
$\displaystyle\lim_{x\to a}\frac{f(x)-f(a)}{x-a}$.\\
Enfin du moins quand c’est réel.

\wcSpan*[class=tp-red-style]{— Euh…}

Tu comprendras quand tu seras plus grande…

\wcSpan*[class=tp-red-style]{— Mais en vrai, c’est quoi une dérivée ?}

\wcTpFlipCard[style=--tp-flip-card-width:26rem;--tp-flip-card-aspect-ratio:0.75;]{
  \begin{description}
    \item[recto]
      Quel est l'\textbf{accroissement} de la fonction $f(x) = x^2$
      entre $x = 2$ et $x = 5$ ?

      \wcTpXyPlot{
        \rawhtml{<script type="tp/xy-plot">
xyFunctionGraph
  x-axis [-1,6]
  x-ticks [-1,0,1,2,3,4,5,6]
  legend-x "x"
  y-ticks [-5,0,1,4,9,16,25,36]
  y-axis [-5,36]
  legend-y "f(x)"
  points [["A",2,4,-12,0], ["B",5,25,-12,0]]
  functions [["Parabola", "x*x"]]
  vectors [["Dx",2,4,3,0,"dashed"], ["Df",5,4,0,21,"dashed"]]
</script>}
      }

    \item[verso]
      \wcTpStack{
        \wcP{
          L’\textbf{accroissement} $\Delta f$ de la fonction $f(x)$ est la variation
        de sa valeur entre $x = a$ et $x = b$ :
        }

        \wcP{
          $\Delta f = f(b) - f(a)$.
        }

        \wcP{
          Avec $f(x) = x^2$, $a = 2$ et $b = 5$, on a :
        }

        \wcP{
          $\Delta f = f(5) - f(2) = 5^2 - 2^2 = 25 - 4 = 21$
        }

        \wcTpBox{
          $\Delta f = 21$ est l'accroissement de $f(x)=x^2$
          entre $x = 2$ et $x = 5$.
        }
      }
  \end{description}
}

\wcTpFlipCard[style=--tp-flip-card-width:26rem;--tp-flip-card-aspect-ratio:0.75;]{
  \begin{description}
    \item[recto]
      Quel est le \textbf{taux d’accroissement} de la fonction $f(x) = x^2$
      entre $x = 2$ et $x = 5$ ?

      \wcTpXyPlot{
        \rawhtml{<script type="tp/xy-plot">
xyFunctionGraph
  x-axis [-1,6]
  x-ticks [-1,0,1,2,3,4,5,6]
  legend-x "x"
  y-ticks [-5,0,1,4,9,16,25,36]
  y-axis [-5,36]
  legend-y "f(x)"
  points [["A",2,4,-12,0], ["B",5,25,-12,0]]
  functions [["Parabola", "x*x"],["Secant", "7*x - 10", [2,5]]]
  vectors [["Dx",2,4,3,0,"dashed"], ["Df",5,4,0,21,"dashed"]]
</script>}
      }

    \item[verso]
      \wcTpStack{
        \wcP{
          Le \textbf{taux d’accroissement} de la fonction $f(x)$ rapporte
        son accroissement $\Delta f$ à la variation correspondante $\Delta x$ de $x$ :
        }

        \wcP{
          $\displaystyle\frac{\Delta f}{\Delta x} = \frac{f(b) - f(a)}{b - a}$
        avec $a \neq b$.
        }

        \wcP{
          Avec $f(x) = x^2$, $a = 2$ et $b = 5$, on a :
        }

        \wcP{
          $\displaystyle\frac{\Delta f}{\Delta x} = \frac{f(5) - f(2)}{5 - 2} = \frac{21}{3} = 7$
        }

        \wcTpBox{
          \begin{itemize}
            \item $\Delta x = 3$ est la variation de $x$.
            \item $\Delta f = 21$ est l'accroissement correspondant de $f(x)=x^2$.
            \item $\displaystyle\frac{\Delta f}{\Delta x} = 7$ est le taux
              d'accroissement de $f(x)=x^2$ entre $x = 2$ et $x = 5$.
          \end{itemize}
        }

        \wcP{
          Géométriquement, le taux d’accroissement est la pente de la droite
        sécante passant par les points $A(a,f(a))$ et $B(b,f(b))$.
        }
      }
  \end{description}
}

\end{document}
